\section{LaTeX, LuaLaTeX ויצירת PDF מקצועי}

\LaTeX{} הוא שפת סימון לייצור מסמכים מקצועיים \cite{lamport1994latex}. במסגרת
ההרצאה הוצג ככלי מרכזי לעידן הסוכנים: מסמך שנכתב בטקסט בלבד, ללא גרירת
עכבר, מתאים להפעלה אוטומטית על ידי \eng{AI}.

\subsection{מדוע LaTeX?}

\begin{itemize}[rightmargin=2em]
  \item \textbf{רפטביליות} --- אותו קוד מייצר אותו \eng{PDF} בכל הרצה.
  \item \textbf{הפרדה} --- תוכן (\texttt{.tex}) נפרד מעיצוב (\texttt{.cls}).
  \item \textbf{מסמכים גדולים} --- תוכן עניינים, ביבליוגרפיה, הפניות אוטומטיות.
  \item \textbf{מתמטיקה} --- נוסחאות ברמה ש-\eng{Word} מתקשה להשיג.
  \item \textbf{סוכן-ידידותי} --- פקודות טקסט במקום קואורדינטות עכבר.
\end{itemize}

המרצה הצהיר במפורש: ``טעות פטאלית לעבוד בוורד'' למסמכים רציניים וגדולים.

\subsection{סוגי קבצי מסמך}

ההרצאה סקרה משפחת מסמכי טקסט:
\begin{description}[rightmargin=2em, style=nextline]
  \item[\texttt{.txt}] טקסט שטוח --- ללא עיצוב, מתאים לקוד וסקילים.
  \item[\texttt{.md}] \eng{Markdown} --- כותרות, הדגשות, תמונות; פופולרי לתיעוד.
  \item[\texttt{.tex}] \LaTeX{} --- מסמכים אקדמיים ומקצועיים ל-\eng{PDF}.
  \item[\texttt{.html}] דפי אינטרנט --- ניתן להמיר ל-\LaTeX{} על ידי סוכן.
\end{description}

\subsection{LuaLaTeX --- הקומפיילר}

\eng{LuaLaTeX} הוא מנוע קומפילציה מבוסס \eng{Lua}, התומך ב-\eng{Unicode}
ובפונטים מודרניים (דרך \eng{fontspec}) \cite{polyglossia2023}. לעברית
זה קריטי: \eng{pdfLaTeX} אינו מספיק לשילוב עברית-אנגלית איכותי.

תהליך הקומפילציה:
\begin{enumerate}[rightmargin=2em]
  \item כותבים \texttt{main.tex} + \texttt{.cls} + \texttt{.bib}.
  \item מריצים \texttt{lualatex main.tex}.
  \item מריצים \texttt{bibtex main} (אם יש ביבליוגרפיה).
  \item מריצים \texttt{lualatex} פעמיים נוספות (תוכן עניינים והפניות).
\end{enumerate}

\begin{lstlisting}[caption={דוגמת שימוש בפקודות התבנית}]
\documentclass{agentic-ai-template}
\begin{document}
\maketitlepage
שילוב עברית עם \eng{API} ו-\eng{RAG}.
\begin{notebox}
תיבת מידע מעוצבת מהקובץ \texttt{.cls}
\end{notebox}
\end{document}
\end{lstlisting}

\subsection{מבנה פרויקט לסוכן}

לפי דוגמת המרצה:
\begin{itemize}[rightmargin=2em]
  \item \texttt{main.tex} --- קובץ ראשי, יכול לכלול \texttt{\textbackslash input\{\}}.
  \item \texttt{agentic-ai-template.cls} --- כל העיצוב והפקודות המותאמות.
  \item \texttt{references.bib} --- מקורות בפורמט \eng{BibTeX}.
  \item \texttt{sources/} --- חומר גלם (תמלול, סיכומים).
\end{itemize}

\subsection{יתרונות לעבודה עם סוכן AI}

\begin{insightbox}
סוכן שכותב \LaTeX{} מייצר \eng{diff} קריא, ניתן לגרסה ב-\eng{Git}, וניתן
לתיקון נקודתי. מסמך \eng{Word} שנוצר על ידי סוכן לעיתים ``שובר'' עיצוב
בלתי נראה עד להדפסה.
\end{insightbox}

\subsection{מגבלות}

\begin{itemize}[rightmargin=2em]
  \item עקומת למידה --- שגיאות קומפילציה דורשות הבנה בסיסית.
  \item עברית-\eng{BiDi} --- בעיות RTL/LTR דורשות תבנית מותאמת.
  \item התקנת חבילות --- \eng{MiKTeX}/\eng{TeX Live} נדרשים מקומית.
\end{itemize}

פרקים 10--11 יראו כיצד \eng{CLS} וסקילי \eng{QA} פותרים חלק מהמגבלות
הללו בצורה רפטבילית.
